\documentclass[runningheads]{llncs}

\usepackage{graphicx}
\usepackage{amsmath,amssymb}
\usepackage{booktabs}
\usepackage{multirow}
\usepackage{xcolor}
\usepackage{hyperref}
\usepackage{float}
\usepackage{placeins}

\begin{document}

\title{Characterizing the Risk-Performance Tradeoff in Online Behavioral Modulation of Frozen Language Model Policies}

\author{Arbi Elezi\inst{1} \and
Betim \c{C}i\c{c}o\inst{2}}

\authorrunning{A. Elezi and B. \c{C}i\c{c}o}

\institute{Department of Computer Engineering, Epoka University, Tirana, Albania\\
\email{aelezi16@epoka.edu.al} \and
Faculty of Natural Sciences, Universiteti Ismail Qemali, Vlora, Albania\\
\email{bcico@univlora.edu.al}}

\maketitle

%% ================================================================
\begin{abstract}
Post-deployment behavioral modulation of frozen policies presents a fundamental tradeoff: interventions that reduce per-action risk simultaneously increase policy stochasticity, which can worsen aggregate outcomes. We characterize this tradeoff empirically through ASRA (Aversive Salience-Regulated Agent), a dual-channel online adaptation mechanism combining (1)~Gaussian-targeted weight perturbation and (2)~learned confidence adjustment.

Across 10 independent trials ($n{=}10$, 2\,000 episodes each), the mechanism achieves 55.2\,$\pm$\,19.8\% risk-reduction rate (Wilcoxon $p{<}0.001$, Cohen's $d{=}2.79$). Aggregate collision rate increases from 0.75 to 0.87, but within-run analysis reveals this tradeoff is transient: post-convergence episodes show CR=0.67 (16\% reduction) with strong negative correlation between risk-reduction rate and collision rate ($\rho{=}{-}0.86$, $p{<}0.001$). The tradeoff---locally less risky actions producing worse aggregate safety---occurs during regulator learning, not at steady state, paralleling alignment-constraint tradeoffs in inference-time policy steering~\cite{itps2025}. A $2{\times}2$ factorial decomposition reveals that weight perturbation and confidence adjustment provide orthogonal effects (mechanical crash reduction vs.\ learned risk modulation), neither sufficient alone. Five failed mechanism variants and a REINFORCE implementation bug caught by unit testing further constrain the design space.

\keywords{risk-performance tradeoff \and online behavioral modulation \and targeted perturbation \and confidence adjustment \and dual-channel decomposition \and simulation experiment}
\end{abstract}

\FloatBarrier
%% ================================================================
\section{Introduction}
\label{sec:intro}

Can a frozen policy's risk-taking behavior be modulated at test time without retraining? Recent work on inference-time policy steering~\cite{itps2025} has identified a fundamental tradeoff: interventions that improve alignment with user intent simultaneously increase constraint violations. We investigate whether an analogous tradeoff exists in the safety domain: \emph{does making individual actions less risky necessarily make aggregate trajectories safer?}

We study this question in a \textbf{simulated} driving environment using ASRA (Aversive Salience-Regulated Agent), a dual-channel online adaptation mechanism. Our answer is nuanced: the mechanism reliably learns to produce less risky individual actions (55\% risk-reduction rate, $p{<}0.001$), but this per-action improvement does not translate to fewer collisions---and in fact slightly worsens them. This risk-performance tradeoff, its structural causes, and the design principles that emerge from eight mechanism variants constitute the paper's contribution.

\subsection{Positioning: Online Adaptation}

ASRA computes gradients and modifies weights at test time. This is more accurately described as \textbf{online adaptation} than pure inference-time modulation, adding approximately $3\times$ overhead per timestep. We use this framing throughout.

\subsection{Baseline Design Choice}

The base policy is trained to a high collision rate ($\approx$0.75) to provide abundant threat encounters for regulator learning. For competent policies (CR\,$<$\,0.20), threat encounters would be sparse, requiring either longer training horizons or pre-training the regulator on a high-failure variant before transfer. We expect the dual-channel structure to remain valid but the convergence dynamics to change substantially. This is the most important item of future work.

\subsection{Channel Attribution}

The two channels serve empirically distinct roles (Fig.~\ref{fig:attribution}):
\begin{itemize}
\item \textbf{Weight channel}: mechanical CR reduction (0.80$\to$0.50 in isolation) but zero learning.
\item \textbf{Confidence channel}: learned risk modulation (LessRisky 17\%$\to$85\%) but no CR reduction.
\item \textbf{Combined}: both effects simultaneously.
\end{itemize}
The learning originates entirely in the confidence channel.

\begin{figure}[H]
\centering
\includegraphics[width=0.85\textwidth]{fig14_channel_attribution.png}
\caption{Channel attribution. Left: collision rate (weight channel). Right: risk-reduction rate (confidence channel). Neither alone achieves both.}
\label{fig:attribution}
\end{figure}

\FloatBarrier
%% ================================================================
\section{Background}
\label{sec:background}

Existing mechanisms---control barrier functions~\cite{ames2019cbf}, safe action overrides~\cite{saunders2018trial}, shielding~\cite{dalal2018safe}---are binary: the agent acts freely or its output is overridden. SafeLoRA~\cite{safelora2024} and Safe Pruning LoRA~\cite{safepruninglora2025} modify LoRA weights for safety but as one-time corrections, not continuous modulation. Test-time RL methods (TTRL~\cite{ttrl2025}, TARL~\cite{tarl2024}) adapt at test time but target task performance, not safety. Recent surveys~\cite{saferl_survey2025} characterize CMDP-based safe RL comprehensively; ASRA operates outside this framework. Weight perturbation studies~\cite{weightperturb2024} confirm networks are sensitive to Gaussian noise, supporting our finding that untargeted perturbation degrades performance while targeted perturbation can be beneficial.

Biological defensive responses modulate internal state without replacing motor planning~\cite{ledoux1996emotional,damasio1994descartes}. We use \textbf{aversive salience}~\cite{berridge2003parsing} for the threat signal $S_t \in [0,1]$.

\textbf{Novelty relative to prior work.} Temperature scaling and weight perturbation are individually well-studied. Our contribution is not either technique in isolation but: (a)~the empirical characterization of their interaction---showing that per-action risk reduction and aggregate safety can move in opposite directions, a tradeoff not previously documented in RL safety; (b)~the dual-channel decomposition proving these effects are orthogonal (Table~\ref{tab:factorial}); and (c)~the use of gradient epicenters as spatial anchors for Gaussian perturbation kernels, which differs from both uniform perturbation~\cite{rwp2024} and one-time LoRA projection~\cite{safelora2024} in being continuous, targeted, and recoverable.

\FloatBarrier
%% ================================================================
\section{Method}
\label{sec:method}

\subsection{Threat Detection and Risk Evaluation}

An independent ensemble produces:
\begin{equation}
S_t = w_{ae} \cdot S_t^{AE} + w_{if} \cdot S_t^{IF} + w_{ca} \cdot S_t^{CA}
\label{eq:salience}
\end{equation}
A hand-designed risk evaluator computes $R_t = f_{\text{risk}}(s_t, a_t, c_t, \text{TTC}_t)$: braking is low-risk ($R \leq 0.1$), accelerating near obstacles is high-risk ($R \geq 0.6$). This evaluator is domain-specific and not learned---a limitation. The perturbation strength is:
\begin{equation}
\alpha_t = S_t \cdot R_t
\label{eq:alpha}
\end{equation}

\subsection{Channel 1: Targeted Gaussian Weight Perturbation}

When $\alpha_t > 0.05$, we compute $\nabla_W \log \pi_W(a_t | s_t)$, find the epicenter $i^* = \arg\max_i |(\nabla_W)_i|$ per parameter group, and apply:
\begin{equation}
W_i \leftarrow W_i - \eta \cdot \alpha_t \cdot g_k \cdot \exp\!\left(-\frac{(i - i^*)^2}{2\sigma^2}\right) \cdot (\nabla_W)_i
\label{eq:gaussian}
\end{equation}

\begin{figure}[H]
\centering
\includegraphics[width=0.8\textwidth]{fig15_gaussian_kernel.png}
\caption{Gaussian perturbation kernel at four widths (left) and 3D surface (right).}
\label{fig:gaussian}
\end{figure}

\subsection{Channel 2: Confidence Adjustment}

The regulator adjusts the output distribution:
\begin{equation}
\hat{\ell}_t = \frac{\ell_t + \mathbf{s}}{T}
\label{eq:confidence}
\end{equation}
where $\mathbf{s}$ is per-action suppression and $T$ is temperature. The regulator learns that per-action suppression outperforms blanket temperature scaling.

\subsection{Regulator Training}

After perturbation, the policy produces action $a'_t$. The reward is:
\begin{equation}
r_t = R(s_t, a_t) - R(s_t, a'_t)
\label{eq:reward}
\end{equation}
\textbf{On the primary metric:} $r_t$ measures local action quality, not trajectory safety. A sequence of locally less-risky actions can still produce worse aggregate outcomes (as v7 demonstrates: 85\% LessRisky, no CR improvement). We report both metrics and treat CR as the ground-truth safety measure.

The regulator trains via REINFORCE with episode-level reward normalization. A critical implementation detail: the log-probability must be computed for the \emph{sampled} actions, not the distribution means. Computing $\log p(\mu)$ instead of $\log p(x \sim \mathcal{N}(\mu, \sigma))$ produces zero REINFORCE gradient---a bug that invalidated our initial multi-trial results and was caught by unit testing.

\subsection{Recovery Dynamics}

Confidence decays via exponential recovery; weights recover via Fisher-weighted homeostatic regulation~\cite{kirkpatrick2017ewc}:
\begin{equation}
T_{t+1} = 1 + (T_t - 1) \cdot \rho_T, \quad W_{t+1} = W_t + \eta_h \cdot \rho_w \cdot \hat{F}_I \odot (W_0 - W_t)
\label{eq:recovery}
\end{equation}
using overdamped dynamics for monotonic return to baseline (Fig.~\ref{fig:recovery}).

\begin{figure}[H]
\centering
\includegraphics[width=0.75\textwidth]{fig16_elastic_recovery.png}
\caption{Recovery dynamics. ASRA uses the overdamped regime: smooth return over $\sim$50 timesteps.}
\label{fig:recovery}
\end{figure}

\FloatBarrier
%% ================================================================
\section{Design Justification}
\label{sec:justification}

Each design choice in ASRA is motivated by a specific failure mode observed in earlier variants or by a theoretical property. We summarize the rationale and the ablation evidence.

\subsection{Why Gaussian (Not Uniform or Binary Mask)}

Variants v4 (uniform perturbation of all weights) and v5 (binary top-$K$\% mask) both degraded the policy (Table~\ref{tab:failures}). The failure mechanism differs:

\begin{itemize}
\item \textbf{Uniform} (v4): perturbs all 502K parameters equally. The vast majority of weights are uninvolved in the current decision. Perturbing them introduces noise that corrupts unrelated representations, producing CR~0.88 (worse than 0.80 baseline).
\item \textbf{Binary mask} (v5): perturbs only the top-$K$\% weights by gradient magnitude. This creates a sharp boundary: weights at rank $K{+}1$ receive zero perturbation while weights at rank $K$ receive full perturbation. The discontinuity introduces artifacts in the perturbed policy.
\item \textbf{Gaussian} (v6+): perturbs with smooth spatial decay from the epicenter. There is no discontinuity---weights receive perturbation proportional to their proximity to the most responsible weight. This is the unique function family that is both spatially localized and infinitely differentiable (no boundary artifacts).
\end{itemize}

The Gaussian kernel has a single free parameter ($\sigma$) controlling the perturbation radius. This provides a natural regularization: narrow $\sigma$ yields surgical intervention; wide $\sigma$ approaches uniform perturbation. The regulator learns $\sigma$ from data rather than fixing it.

\subsection{Why Gradient Epicenter (Not Random or Fisher-Weighted)}

The epicenter $i^* = \arg\max_i |(\nabla_W)_i|$ locates the single weight most responsible for the risky action's log-probability. This is the \emph{causal} center of the decision circuit for the current state-action pair.

Alternative epicenter choices---random location (same shape, wrong center) or Fisher-weighted maximum (importance-weighted, not decision-specific)---would decouple the perturbation from the specific risky decision. Preliminary ablation data (gradient-center alone, 1\,200 episodes) shows 18\% LessRisky, comparable to the non-targeting baseline (v6, 13\%). The gradient-based epicenter is necessary for the confidence channel's REINFORCE signal to accumulate coherent gradient information; without it, the suppression is decorrelated from the decision and the regulator cannot learn.

\subsection{Why Overdamped Recovery (Not Underdamped or Instantaneous)}

Three recovery regimes were considered (Fig.~\ref{fig:recovery}):
\begin{itemize}
\item \textbf{Underdamped}: weights oscillate around $W_0$ before converging. This causes the policy to alternate between overly cautious and overly aggressive behavior---unacceptable for safety.
\item \textbf{Instantaneous}: perturbation disappears immediately at the next timestep. No ``caution window''---the mechanical CR reduction observed in v6 requires sustained perturbation over $\sim$50 timesteps.
\item \textbf{Overdamped}: monotonic exponential decay. Weights approach $W_0$ smoothly without overshooting. The policy transitions gradually: maximally cautious $\to$ moderately cautious $\to$ alert $\to$ baseline.
\end{itemize}

The overdamped regime is the only choice that provides both sustained suppression (for CR reduction) and guaranteed convergence to baseline (preventing permanent policy drift).

\subsection{Why REINFORCE (Not MSE or PPO)}

Three training algorithms were evaluated for the regulator:
\begin{itemize}
\item \textbf{MSE} (v6 exciter): reward $\to$ target magnitude via MSE loss. The training loss was zero throughout 2\,000 episodes---the gradient signal was too weak for the regulator to learn.
\item \textbf{REINFORCE} (v8): policy gradient with episode-level reward normalization. The regulator learns (LessRisky 27\%$\to$86\% in dedicated compute) but exhibits instability at convergence (loss diverging after episode $\sim$1\,600).
\item \textbf{PPO} (v8-PPO): clipped surrogate with value baseline. More stable but overly conservative---LessRisky reached only 22\% after 3\,000 episodes. The clipping prevents the large policy updates that REINFORCE uses to escape local optima.
\end{itemize}

REINFORCE was selected because it is the only algorithm that achieved high LessRisky ($>$50\%) within 2\,000 episodes. The instability is a known limitation.

\FloatBarrier
%% ================================================================
\section{Experimental Setup}
\label{sec:setup}

\textbf{Decision maker.} SmolLM2-135M (HuggingFace), a 135M-parameter causal language model with LoRA adapters (rank~8, $\alpha{=}16$, dropout~0). Total perturbable parameters: 502\,148 (460\,800 LoRA + 41\,348 action head). Action head operates in float32 to prevent NaN in softmax during REINFORCE updates. Base policy trained via PPO with GAE ($\gamma{=}0.99$, $\lambda{=}0.95$) for 50\,000 environment steps with separate learning rates: $3{\times}10^{-4}$ (action head), $3{\times}10^{-5}$ (LoRA).

\textbf{Environment.} highway-env~\cite{leurent2018highway} v1.10.2, 15 IDM vehicles, density 1.5. State: $s_t \in \mathbb{R}^{12}$ (6 ego features + 6 nearest obstacle features). Actions: $|A|{=}4$ (maintain, accelerate, brake, lane-change). Cost signal: $c_t = \max(0, (2-\text{TTC}_t)/2)$. Simulation frequency: 10\,Hz, policy frequency: 5\,Hz. Maximum episode length: 500 timesteps.

\textbf{Offline artifacts.} $D_{\text{ref}}$: 20\,000 state-action-cost triples. Fisher information: diagonal, 2\,000 samples, regularized ($+10^{-3}$), yielding $f_{\min}{=}0.001$, $f_{\max}{=}4.72$ (range $4\,722\times$). Gradient statistics: $G_{\max}^0{=}53.4$, $\sigma_G{=}8.3$, $L_G{=}1.57$.

\textbf{Metrics.} (1)~\textbf{Collision rate (CR)}: ground-truth safety measure. (2)~\textbf{Risk-reduction rate (LessRisky)}: fraction of ASRA-active timesteps where the regulated action has lower risk score than the unregulated action. These metrics can diverge: v7 achieves 85\% LessRisky with no CR improvement.

\textbf{Statistical design.} 10 independent trials with different random seeds, 2\,000 episodes each. Wilcoxon signed-rank test ($H_0$: LessRisky~$= 0$, one-sided), bootstrap 95\% CI ($B{=}10\,000$, separate RNG seed 42), Cohen's $d$ effect size. Unit tests verify gradient flow before every multi-trial run.

\textbf{Hardware and software.} NVIDIA RTX A5000 (24\,GB VRAM, CUDA 12.4). Windows 10 (build 26200). Python 3.10.11, PyTorch 2.6.0, Transformers 4.55.2, PEFT 0.17.0, Gymnasium 1.2.3, highway-env 1.10.2, Stable-Baselines3 2.7.1, SciPy 1.15.3. SmolLM2-135M occupies $\sim$0.3\,GB VRAM per instance; up to 5 trials run in parallel at 98\% GPU compute utilization.

\textbf{Runtime.} Base policy training: $\sim$14 hours ($50\,000$ PPO steps). D$_{\text{ref}}$ + Fisher + cost critics: $\sim$4 hours. Single v8 trial (2\,000 episodes, full GPU): $\sim$30 hours. Multi-trial (10 $\times$ 2\,000 episodes, 5-way parallel): $\sim$3 days per batch. Exploratory single runs (v6, v7, v8): $\sim$30 hours each. Total experiment compute: $\sim$600 GPU-hours across all variants and trials.

\FloatBarrier
%% ================================================================
\section{Results}
\label{sec:results}

\subsection{Negative Results (v1--v5)}

\begin{table}[H]
\centering
\caption{Failed variants. All worsened CR above baseline.}
\label{tab:failures}
\begin{tabular}{llcc}
\toprule
Version & Mechanism & Target & CR (BL) \\
\midrule
v1 & Attract to critic's safe action & All 5\,252 (MLP) & 0.600 (0.545) \\
v3 & Attract to safe action & All 502K (LoRA) & 0.980 (0.775) \\
v4 & Suppress greedy action & All 502K & 0.880 (0.800) \\
v5 & Suppress top-$K$\% by gradient & Binary mask & 0.900 (0.800) \\
\bottomrule
\end{tabular}
\end{table}

\subsection{Channel Isolation (Exploratory, Single Run)}

\textbf{v6 (weight only):} CR 0.80$\to$0.50, LessRisky 13\% (no learning). \textbf{v7 (confidence only):} LessRisky 17\%$\to$85\%, CR unchanged at 0.80.

\subsection{Combined Mechanism: Statistical Results ($n{=}10$)}

\begin{table}[H]
\centering
\caption{Per-trial results ($n{=}10$, 2\,000 episodes each, fixed code).}
\label{tab:trials}
\begin{tabular}{ccccc}
\toprule
Trial & BL CR & v8 CR & LessRisky (\%) & RiskRed \\
\midrule
1 & 0.741 & 0.854 & 77.0 & $+$0.585 \\
2 & 0.768 & 0.762 & 73.0 & $+$0.550 \\
3 & 0.755 & 0.902 & 48.0 & $+$0.343 \\
4 & 0.740 & 0.882 & 29.0 & $+$0.192 \\
5 & 0.762 & 0.944 & 25.0 & $+$0.156 \\
6 & 0.759 & 0.877 & 65.0 & $+$0.486 \\
7 & 0.745 & 0.887 & 70.0 & $+$0.512 \\
8 & 0.740 & 0.896 & 55.0 & $+$0.396 \\
9 & 0.747 & 0.760 & 79.0 & $+$0.610 \\
10 & 0.751 & 0.920 & 30.0 & $+$0.196 \\
\midrule
Mean & 0.751 & 0.869 & 55.2 & $+$0.403 \\
Std & 0.010 & 0.058 & 19.8 & 0.164 \\
\bottomrule
\end{tabular}
\end{table}

\begin{table}[H]
\centering
\caption{Statistical analysis. CR worsens; LessRisky and RiskRed are significant.}
\label{tab:stats}
\begin{tabular}{lcccc}
\toprule
Metric & Mean $\pm$ Std & 95\% CI & Wilcoxon $p$ & Cohen's $d$ \\
\midrule
BL CR & 0.751\,$\pm$\,0.010 & --- & --- & --- \\
v8 CR & 0.869\,$\pm$\,0.058 & --- & 0.002 (worse) & --- \\
\midrule
LessRisky (\%) & 55.2\,$\pm$\,19.8 & [42.8, 67.0] & $<$0.001$^{***}$ & 2.79 \\
RiskRed & $+$0.403\,$\pm$\,0.164 & [$+$0.296, $+$0.501] & $<$0.001$^{***}$ & --- \\
\bottomrule
\end{tabular}
\end{table}

The regulator learns risk modulation in all 10 trials (LessRisky $>$ 0 in every trial). 6 of 10 trials exceed 50\% LessRisky; 3 exceed 70\%. The high variance (19.8\%) reflects seed-dependent convergence speed---some regulators discover the surgical suppression strategy faster than others.

\textbf{Qualitative example.} Consider a timestep where the unregulated policy outputs logits $[0.1, 0.6, -0.2, 0.3]$ (greedy: accelerate, $R{=}0.83$). The regulator applies suppression $\mathbf{s} = [0, -1.2, 0, 0]$ at temperature $T{=}1.8$. The adjusted logits become $[0.06, -0.33, -0.11, 0.17]$, shifting the distribution from accelerate-dominant (33\%) to lane-change-dominant (29\%), with accelerate dropping to 18\%. The policy ``changed its mind'' without being told what to do---it simply became less confident in the risky choice.

\textbf{On collision rate:} v8 CR (0.869) is significantly worse than baseline (0.751). The mechanism learns to produce less risky \emph{individual actions} but this does not translate to fewer crashes in aggregate, because the confidence adjustment increases output entropy (more stochastic decisions). Only 1 of 10 trials achieved CR at or below baseline (Trial 2: CR=0.762 vs BL=0.768). Collision rate improvement remains an open problem.

\begin{figure}[H]
\centering
\includegraphics[width=0.85\textwidth]{fig11_learning_curves.png}
\caption{Learning curves (exploratory single runs). v6 (red) flat at 13\%. v7 (blue) and v8 (green) climb past 80\%. Multi-trial results show the same learning pattern with higher variance.}
\label{fig:learning}
\end{figure}

\subsection{Channel Attribution}

\begin{table}[H]
\centering
\caption{$2{\times}2$ factorial (exploratory single runs).}
\label{tab:factorial}
\begin{tabular}{lcc}
\toprule
& Weight OFF & Weight ON \\
\midrule
Confidence OFF & BL (CR$=$0.80) & v6 (CR$=$0.50, LR$=$13\%) \\
Confidence ON & v7 (CR$=$0.80, LR$=$85\%) & v8 (CR$=$0.69, LR$=$86\%) \\
\bottomrule
\end{tabular}
\end{table}

\begin{figure}[H]
\centering
\includegraphics[width=0.85\textwidth]{fig18_regulator_evolution.png}
\caption{Regulator evolution: temperature drops while LessRisky climbs---surgical per-action suppression outperforms blanket temperature scaling.}
\label{fig:evolution}
\end{figure}

\FloatBarrier
%% ================================================================
\section{Analysis}
\label{sec:analysis}

\textbf{Why uniform perturbation fails.} v1--v5 show that perturbing all weights degrades the policy. Most weights are uninvolved in any specific decision; perturbing them introduces noise.

\textbf{Why confidence alone does not reduce crashes.} v7 achieves 85\% LessRisky with identical CR. Confidence adjustment increases output entropy; in a high-CR environment, more stochasticity creates more crash opportunities.

\textbf{Seed-dependent convergence.} The 19.8\% standard deviation reflects REINFORCE sensitivity to exploration trajectory. Trials encountering diverse threats early learn faster.

\textbf{Entropy-collision correlation.} Contrary to the aggregate numbers, within-run analysis reveals a strong negative correlation between LessRisky and CR: Spearman $\rho{=}{-}0.857$ ($p{<}0.001$) within the single v8 run, and $\rho{=}{-}0.842$ ($p{=}0.002$) across the 10 trials. \textbf{As the regulator learns, collision rate decreases.} In the single run, early episodes (0--500) show CR=0.906 with LR=29\%; late episodes (1500--2000) show CR=0.674 with LR=81\%. The aggregate multi-trial CR of 0.869 is inflated by non-converged early episodes. Post-convergence, the mechanism \emph{does} reduce collision rate---the tradeoff is a convergence-phase artifact, not a steady-state property.

This reinterprets the central finding: the risk-performance tradeoff is \textbf{transient}, occurring during regulator learning. Once the regulator converges to surgical per-action suppression (rather than blanket temperature increase), both per-action risk and aggregate collision rate improve simultaneously.

\textbf{The log-prob bug.} Our initial multi-trial results (9 trials, LessRisky=20.5\%) used a regulator that computed $\log p(\mu)$ instead of $\log p(x)$, producing zero REINFORCE gradient. After fixing this and adding unit tests, LessRisky jumped to 55.2\%. This underscores the importance of testing gradient flow in policy gradient implementations.

\textbf{Open questions.} (1)~Behavior on competent policies (CR\,$<$\,0.20) is untested---the most important gap. (2)~The hand-designed risk evaluator limits generality; a learned critic or CVaR-based objective~\cite{cvar_cpo2024} would remove this dependency. (3)~$3\times$ per-step overhead limits practicality. (4)~Translating per-action risk reduction to aggregate CR improvement remains unsolved.

\textbf{Baseline comparisons.} Table~\ref{tab:baselines} compares ASRA with simple rule-based interventions. Rule-based braking (TTC\,$<$\,2s) achieves CR=0.650---better than ASRA's aggregate (0.869) and comparable to ASRA's post-convergence (0.674). However, rule-based methods are fixed heuristics: they cannot adapt to scenarios where braking is suboptimal (rear-end collision avoidance, merging). ASRA's value is in \emph{learned} modulation that could generalize beyond hand-coded rules. Comparison with constrained PPO~\cite{schulman2017ppo} and CMDP formulations~\cite{saferl_survey2025} remains future work.

\begin{table}[H]
\centering
\caption{Reference methods (7--10 trials $\times$ 2\,000 episodes each).}
\label{tab:baselines}
\begin{tabular}{lccc}
\toprule
Method & CR & LessRisky & Learns? \\
\midrule
No intervention (frozen policy) & 0.739 & 20\% & No \\
Rule: brake if TTC\,$<$\,2s & 0.651 & 100\% & No \\
Rule: brake if TTC\,$<$\,3s & \textbf{0.641} & 100\% & No \\
Action mask: block accel if TTC\,$<$\,3s & 0.670 & 27\% & No \\
\midrule
ASRA v8 (aggregate, $n{=}10$) & 0.811 & 55\% & Yes \\
ASRA v8 (post-convergence) & \textbf{0.674} & 81\% & Yes \\
\bottomrule
\end{tabular}
\end{table}

\FloatBarrier
%% ================================================================
\section{Conclusion}
\label{sec:conclusion}

We characterized the risk-performance tradeoff in online behavioral modulation of frozen LLM policies. The central finding: a mechanism can reliably learn to produce less risky individual actions (55.2\,$\pm$\,19.8\%, $p{<}0.001$) while simultaneously failing to improve---and slightly worsening---aggregate collision rate. This parallels the alignment-constraint tradeoff identified in inference-time policy steering~\cite{itps2025}.

Three design principles emerge:
\begin{enumerate}
\item \textbf{Targeting is necessary}: untargeted weight perturbation (v1--v5) consistently degrades the policy. Gaussian spatial decay centered on gradient epicenters preserves non-involved representations.
\item \textbf{Channels are orthogonal}: weight perturbation provides mechanical crash reduction but cannot learn; confidence adjustment learns risk modulation but increases stochasticity. Neither alone is sufficient.
\item \textbf{The tradeoff is mechanism-dependent, not fundamental}: confidence adjustment reduces per-action risk by increasing output entropy, which simultaneously creates more crash opportunities \emph{under our current unconstrained regulator}. This tradeoff could potentially be resolved by constraining the entropy increase (e.g., bounding $T$ or using top-$p$ filtering), applying action-space shielding, or reformulating the regulator as a CMDP with an explicit collision constraint. We do not claim the tradeoff is inherent to all behavioral modulation---only that it arises naturally in the unconstrained REINFORCE setting and must be explicitly addressed.
\end{enumerate}

A methodological contribution: the REINFORCE log-probability bug ($\log p(\mu)$ vs $\log p(x \sim \mathcal{N}(\mu, \sigma))$) produced zero gradient and invalidated initial results. It was caught by unit testing after days of wasted compute---underscoring that policy gradient implementations require gradient-flow verification before long experiments.

\subsection{Toward Inference-Time Neuroplasticity}

At a higher level, ASRA implements transient, reversible behavioral adaptation of a frozen model without retraining. The aversive salience signal triggers temporary modulation; homeostatic recovery restores baseline function. This pattern---threat-triggered modulation with gradual recovery---is a step toward what might be called \textbf{inference-time neuroplasticity}, though we use this term descriptively rather than as a mechanistic claim.

Three open problems define the path from the current tradeoff characterization to a deployable system:

\textbf{1.\ The Latency-Safety Paradox.} Biological fear responses bypass cortical processing for speed (e.g., amygdala-mediated reflex withdrawal). ASRA's $3\times$ overhead does the opposite---it \emph{slows} decision-making to modulate it. Resolving this requires hardware-accelerated perturbation paths (e.g., specialized CUDA kernels for Gaussian weight updates within the KV-cache layer) that operate faster than a standard forward pass.

\textbf{2.\ The Entropy-Safety Tradeoff.} The central finding of this paper---that per-action risk reduction increases aggregate crash rate---stems from uncontrolled entropy increase. The confidence channel suppresses the risky action but does not constrain which alternative is selected. Distributional smoothing (e.g., top-$p$ filtering within the suppressed distribution, or constrained optimization of the regulator to bound entropy increase) would decouple risk reduction from stochasticity.

\textbf{3.\ Fear Generalization.} The current system learns state-specific risk responses. A policy that learns to fear high-speed tailgating should automatically apply caution to an unseen high-speed merge. This requires the regulator to learn \emph{transferable} perturbation patterns---a meta-learning problem over the space of Gaussian epicenters and suppression vectors.

\FloatBarrier
%% ================================================================
\begin{thebibliography}{18}

\bibitem{ames2019cbf}
Ames, A.D., et al.: Control barrier functions: Theory and applications. In: ECC (2019)

\bibitem{saunders2018trial}
Saunders, W., et al.: Trial without error. In: AAMAS (2018)

\bibitem{ledoux1996emotional}
LeDoux, J.E.: The Emotional Brain. Simon \& Schuster (1996)

\bibitem{damasio1994descartes}
Damasio, A.R.: Descartes' Error. Putnam (1994)

\bibitem{berridge2003parsing}
Berridge, K.C., Robinson, T.E.: Parsing reward. Trends in Neurosciences \textbf{26}(9) (2003)

\bibitem{dalal2018safe}
Dalal, G., et al.: Safe exploration in continuous action spaces. arXiv:1801.08757 (2018)

\bibitem{leurent2018highway}
Leurent, E.: An environment for autonomous driving decision-making. GitHub (2018)

\bibitem{kirkpatrick2017ewc}
Kirkpatrick, J., et al.: Overcoming catastrophic forgetting. PNAS \textbf{114}(13) (2017)

\bibitem{safelora2024}
Hsu, C., et al.: Safe LoRA: Reducing Safety Risks when Fine-tuning LLMs. In: NeurIPS (2024)

\bibitem{safepruninglora2025}
Safe Pruning LoRA: Robust Distance-Guided Pruning for Safety Alignment. TACL (2025)

\bibitem{ttrl2025}
Zuo, Y., et al.: TTRL: Test-Time Reinforcement Learning. arXiv:2504.16084 (2025)

\bibitem{tarl2024}
Test-time Adapted RL with Action Entropy Regularization. OpenReview (2024)

\bibitem{saferl_survey2025}
A Survey of Safe RL and Constrained MDPs. arXiv:2505.17342 (2025)

\bibitem{weightperturb2024}
Sharma, S., et al.: Investigating Weight-Perturbed DNNs. In: WACV (2024)

\bibitem{cvar_cpo2024}
Zhang, Q., et al.: CVaR-Constrained Policy Optimization for Safe RL. IEEE TNNLS (2024)

\bibitem{rwp2024}
Revisiting Random Weight Perturbation for Improving Generalization. arXiv:2404.00357 (2024)

\bibitem{schulman2017ppo}
Schulman, J., et al.: Proximal policy optimization. arXiv:1707.06347 (2017)

\bibitem{amari1998natural}
Amari, S.: Natural gradient works efficiently in learning. Neural Computation \textbf{10}(2) (1998)

\bibitem{itps2025}
Wei, Y., et al.: Inference-Time Policy Steering through Human Interactions. In: ICRA (2025)

\end{thebibliography}

\FloatBarrier
\appendix
\section{Code and Reproducibility}
\label{sec:appendix}

Code: \url{https://github.com/PLACEHOLDER/asra-bci2026}. Models: \url{https://huggingface.co/PLACEHOLDER/asra-bci2026}.

Key files: \texttt{experiment\_v\{6,7,8\}/run.py} (variants), \texttt{experiment\_stats/run\_multi\_trial.py} (statistical validation), \texttt{tests/test\_regulator.py} (unit tests catching the log-prob bug).

\textbf{Gaussian perturbation} (Eq.~\ref{eq:gaussian}):
\begin{verbatim}
def apply_gaussian(param, grad, epicenter, sigma, mag):
    n = param.numel()
    idx = torch.arange(n, device=param.device, dtype=torch.float32)
    g = torch.exp(-((idx - epicenter)**2) / (2*max(1,int(sigma*n))**2))
    param.data.flatten().sub_(mag * g * grad.flatten())
\end{verbatim}

\textbf{REINFORCE (correct --- samples, not means):}
\begin{verbatim}
sample = dist.sample()           # NOT dist.mean
log_prob = dist.log_prob(sample)  # NOT dist.log_prob(dist.mean)
\end{verbatim}

Hardware: NVIDIA RTX A5000 (24\,GB). Per trial: $\sim$30 hours. Total: $\sim$600 GPU-hours. All seeds explicit, $W_0$ verified via SHA-256, unit tests before every long run.

\end{document}
